\subsection{N-014: the printed highest-weight label order is not defined}
\label{sec:nonstandard-label-order}

\paragraph{The precise source rule.}
GCT VIII, subsection 2.3, printed p.~11
(arXiv:0709.0751v2), defines $\mu(\alpha)$ to be the highest
$H_q$-weight in the restriction of $W_{q,\alpha}$ and sets
$\alpha\le\alpha'$ if and only if
$\mu(\alpha)\le\mu(\alpha')$ in the usual root-cone order.
The text permits several highest-weight vectors, but does not
specify a set of incomparable highest weights or a tie rule
for distinct labels with the same highest weight.
Subsection 3.3, step (1), printed p.~20, transfers this order
to the nonstandard Hecke labels by duality; step (2) then
orders GT tableaux lexicographically.

There are two distinct problems with this literal definition
already among the rank-four modules in the supplied examples.
First, the required single greatest weight need not exist.
Second, even where it exists, distinct labels can have the same
weight, so the stated ``if and only if'' gives a preorder,
not a partial order on those labels.  The argument below
uses exact modules, rather than relying on the floating-point
qualification in VII's experimental table.

\paragraph{An actual module with two incomparable maximal weights.}
Let $M_\lambda$ denote the ordinary Hecke simple of shape
$\lambda\vdash4$.  Write $N=M_4\otimes M_{22}$, one of the
two-dimensional simples proved distinct from all other factors
in N-010.  For $H_q=GL_q(2)\times GL_q(2)$, let $V_\lambda$
be the standard quantum Weyl module of shape $\lambda$.
Ordinary quantum Schur--Weyl duality and N-010 give the exact
$H_q$-module identity
\[
 \operatorname{Hom}_{B_{4,K}}(N,X_q^{\otimes4})
       \cong V_4\otimes V_{22}\ \oplus\ V_{22}\otimes V_4.
 \tag{N-53}
\]
Here the Hom denotes the multiplicity space for the actual
nonstandard action on the tensor power, and $K=\mathbb Q(v)$.
In the duality terminology of VII this is $W_2$.

For completeness, the needed exact tensor restrictions for the
three ordinary shapes with at most two rows are
\[
\begin{array}{c|c}
\text{ordinary tensor product}&\text{nonstandard constituents}\\\hline
M_4\otimes M_4&\epsilon_+\\
M_4\otimes M_{31}&M_4\otimes M_{31}\\
M_4\otimes M_{22}&N\\
M_{31}\otimes M_{31}&\epsilon_+\oplus\Lambda^2M_{31}\oplus S'M_{31}\\
M_{31}\otimes M_{22}&M_{31}\otimes M_{22}\\
M_{22}\otimes M_{22}&\epsilon_+\oplus\epsilon_-\oplus S'M_{22}.
\end{array}
 \tag{N-54}
\]
Tensor flip gives the reversed pairs.  The symmetric and
alternating decompositions, the simultaneous $D$-eigenline,
the quotient models and pairwise distinctness of all listed
simples were verified exactly in N-010.  The alternating
square of $M_{22}$ has $P_i=0$: its dimension is one and a
rank-one matrix $C'_i$ has tensor eigenvalues yielding zero
on that alternating line.  Thus $N$ occurs once in each of
$M_4\otimes M_{22}$ and its flip, and nowhere else in (N-54).
This proves (N-53) without a numerical specialization and
without assuming the entire nonstandard duality conjecture.
If that conjecture is used to name the paired quantum-group
module, its restriction must be this multiplicity space.

The two highest weights in (N-53) are
\[
 \lambda=((4,0),(2,2)),\qquad
 \mu=((2,2),(4,0)).
 \tag{N-55}
\]
Their difference is $(2,-2;-2,2)$.  The simple-root cone of
$GL_2\times GL_2$ consists of
$(a,-a;b,-b)$ with $a,b\ge0$.  Neither difference belongs
to this cone.  Each weight in (N-55) is maximal in the
direct sum, and they are incomparable.  Therefore there is
no single highest weight $\mu(N)$ in the sense required
by VIII's displayed order rule.

This is exactly the decomposition displayed in VII,
subsection 7.2.2, printed p.~48 (arXiv:0709.0749v2).
The accompanying computational qualification is on printed
p.~49.  Equations (N-53)--(N-54) supply an exact independent
justification for the particular decomposition used here.

\paragraph{The defect persists for a faithful full-rank realization.}
The preceding example is not repaired by increasing the
dimensions of the two standard vector spaces so that all
ordinary Hecke shapes occur.  To see this, the involution
$\theta(T_i)=-T_i^{-1}$ conjugates ordinary Specht shapes
and twisting one tensor factor sends $P_i$ to $D-P_i$.
Thus the remaining shapes $211$ and $1111$ are obtained
from $31$ and $4$ by the same twist.  The shape $22$ is
self-conjugate.  On the explicit $N$ matrices of N-010,
the intertwiner
\[
       J=\begin{pmatrix}1&1\\-D&-1\end{pmatrix},\qquad
       \det J=D-1,
       \qquad J P_i=(D I-P_i)J
 \tag{N-56}
\]
verifies directly that $N$ is invariant up to isomorphism
under this nonstandard twist.

Twisting both ordinary tensor factors leaves the nonstandard
action unchanged.  Twisting one factor preserves dimensions
and simplicity of constituents.  Consequently the products
involving a shape $31$ or $211$ have constituents of
dimensions $1,3,5$ or the single dimension $6$, as applicable,
and cannot add $N$ to their decompositions.  The product
$M_{22}\otimes M_{22}$ is itself invariant under the twist;
its one-dimensional constituents are interchanged, so its
unique two-dimensional constituent $S'M_{22}$ is invariant
and remains distinct from $N$.
It follows that, among all five shapes of size four, $N$
occurs exactly in the four ordered pairs
\[
          (4,22),\quad(22,4),\quad(1111,22),\quad(22,1111).
 \tag{N-57}
\]
Each occurs with multiplicity one.  For $GL_q(d)\times GL_q(e)$,
the last two summands are omitted when the relevant
four-row Weyl module vanishes.  When present their highest
weights are dominated by those of the corresponding first
two summands.  The maximal antichain remains
$\{(4,22),(22,4)\}$ for every $d,e\ge2$.
Thus the failure of a greatest weight applies also when
the full rank-$114$ algebra acts faithfully.

\paragraph{Equal highest weights are a separate issue.}
For $GL_q(2)\times GL_q(2)$, (N-54) also proves
\[
 W_{\epsilon_-}\cong W_{S'M_{22}}
          \cong V_{22}\otimes V_{22},\qquad
 W_{\Lambda^2M_{31}}\cong W_{S'M_{31}}
          \cong V_{31}\otimes V_{31}
 \tag{N-58}
\]
as $H_q$-modules.  The respective nonstandard simple labels
are distinct by N-010.  VII's notation for these pairs is
$(W_1,W_6)$ and $(W_3,W_7)$, respectively, in subsection 7.2.2.
The literal comparison through their equal highest weights
therefore relates each pair in both directions, violating
antisymmetry on labels.  One may quotient the label set by
this equivalence, but then the inductive subspace attached
to a minimal type need not be a single GT line.  That is a
different prescription, not a consequence of the stated
partial-order definition.

\paragraph{Explicit alternative orders, and the remaining boundary.}
One possible repair assigns to every label its finite
antichain $A_\alpha$ of maximal $H_q$-weights and compares
these antichains by
\[
 A\preceq B\quad\Longleftrightarrow\quad
 \text{for each }a\in A\text{ there exists }b\in B
                       \text{ with }a\le b.
 \tag{N-59}
\]
This is a partial order on antichains: equivalently, it is
inclusion of their generated lower sets.  A partial order on
the original labels is obtained by using strict inclusion
for distinct antichains and declaring distinct labels with
equal antichains incomparable.  A total extension may instead
resolve every tie by an explicitly recorded additional rule.
For example, either $\epsilon_-<S'M_{22}$ or its reverse
is possible after replacing equality of their antichains
by a chosen tie order; likewise either order is possible
for $\Lambda^2M_{31}$ and $S'M_{31}$ in the two-row example.
These choices alter the lower-label spans in the inductive
globalization and therefore cannot be left implicit.

Another possible convention chooses one maximal weight
according to a fixed total refinement of the product root
order, and then separately resolves equal weights.  For
(N-55), giving the first factor priority chooses $\lambda$;
giving the second factor priority chooses $\mu$.  Those are
two different extra choices, neither specified by VIII.

No theorem here says that a repaired order cannot produce a
useful basis.  The conclusion is the precise source-level
one: the rank-four GT data in N-010 do not complete VIII's
printed label-order recipe.  Together with N-012, this
explains why computing a table cannot by itself identify
that table with a uniquely prescribed $E(4)$.
